\documentclass[conference]{IEEEtran}
\usepackage[T1]{fontenc}
\usepackage{amsmath}
\usepackage{graphicx}
\usepackage{float}

\title{Capacitor Charging and a Power On Reset for an Active Low Microcontroller}
\author{
    \IEEEauthorblockN{Kenneth Dialogo}
    \IEEEauthorblockA{Department of Computer Engineering Technology\\
    Vanier College}
    }
\date{}

\begin{document}
\maketitle

\begin{figure}[t]
    \centering
    \includegraphics[width=\columnwidth]
        {rc_charging.pdf}
    \caption{Charging curve for
        $R = 10\,\text{k}\Omega$, $C = 22\,\text{pF}$.}
        \label{fig:rc}
\end{figure}

\begin{abstract}
This document derives the charging behavior of a resistor-capacitor (RC) network, the voltage reached after one time constant, and applies the result to size a power-on-reset (POR) circuit for a microcontroller with an active-low reset.
\end{abstract}

\section{Introduction}
When a capacitor C is charged through a series resistor R from a constant supply $V_0$, the voltage of the capacitor increases according to the spectra of the resistor.
\begin{equation}
V(t) = V_0\left(1-e^{-t/RC}\right),
\label{eq:one}
\end{equation}
where $\tau = RC$ is the time constant. Substituting $t = \tau$ gives $V(\tau) = V_0\left(1-e^{-t/RC}\right) \approx 0.632\,V_0$.

\section{Power On Reset With an Active Low Reset}

A power-on-reset circuit exploits \eqref{eq:one} directly: R from $V_CC$ to the reset node and $C$ from that node to ground hold the microcontroller's active-low reset pin low until the node crosses $V_1H$, releasing the device.

\begin{figure}[H]
    \centering
    \includegraphics[width=0.4\columnwidth]
        {por_circuit.pdf}
    \caption{Power-on-reset network driving an active-low reset pin.}
        \label{fig:VDC}
\end{figure}

\subsection{Sizing R for a Timing Constraint}
The reset pin releases the instant the node voltage reaches $V_{IH}$.
Setting $V(t_{\text{release}}) = V_{IH}$ in \eqref{eq:one} and solving
for $t_{\text{release}}$:
\begin{align*}
V_{IH} &= V_0\left(1-e^{-t_{\text{release}}/RC}\right) \\
\frac{V_{IH}}{V_0} &= 1-e^{-t_{\text{release}}/RC} \\
e^{-t_{\text{release}}/RC} &= 1-\frac{V_{IH}}{V_0} \\
-\frac{t_{\text{release}}}{RC} &= \ln\left(1-\frac{V_{IH}}{V_0}\right)
\end{align*}


\begin{equation}
t_{\text{release}} = -RC\ln\left(1-\frac{V_{IH}}{V_0}\right).
\label{eq:t_release}
\end{equation}

Rearranging \eqref{eq:t_release} for $R$, given a target $C:$
and a chosen $C$:
\begin{equation}
R = \frac{-t_{\text{release}}}{C\ln\left(1-\dfrac{V_{IH}}{V_0}\right)}.
\label{eq:R_sizing}
\end{equation}

\indent 1) \textit{Worked example}: Take $V_{IH} = 0.7\,V_0$, $V_0 = 5.0\,$V, and
a target $t_{\text{release}} = 10\,$ms. The $C = 22\,$pF used for the above plots
 was chosen only to keep the curves readable on a sub-microsecond
timescale; a practical timing design instead starts from $C = 100\,$nF.
Evaluating the denominator of \eqref{eq:R_sizing}:
\begin{equation}
\ln(1-0.7) = \ln(0.3) \approx -1.204.
\label{eq:denom}
\end{equation}

\noindent Substituting into \eqref{eq:R_sizing}:
\begin{equation}
R = \frac{-(10\times10^{-3})}{(100\times10^{-9})\times(-1.204)}
\approx 83\,\text{k}\Omega.
\label{eq:R_result}
\end{equation}

\noindent An $83\,\text{k}\Omega$ resistor with a $100\,$nF capacitor gives a reset hold-off of approximately $10\,$ms. Keeping $C$ at the plotting value of

\section{Applications and Limitations}

\subsection{Application: LED Lamp Appliance Controller}
A plug-in LED lamp with an MCU cycling through
off/dim/bright/fade modes via a button is well served by an
RC-only power-on reset. The product is cost-sensitive and
non-safety-critical, and does not require brown-out
detection, watchdog-triggered resets, or precise timing,
it only needs a clean startup once plugged in. Choosing
$R = 150\,\text{k}\Omega$ and $C = 100\,\text{nF}$ with
$V_0 = 5.0\,\text{V}$ gives, from \eqref{eq:t_release},
\[
t_\text{release} \approx 18.1\,\text{ms}.
\]

\subsection{Limitations}
Three limitations bound the reliability of this design.
First, component tolerance (Section~\ref{sec:tolerance})
causes the true release time to deviate from the nominal
calculation. Second, the circuit provides no brown-out
detection: a brief sag in the wall adapter's output during
operation is invisible to a circuit that only reacts once,
at power-up. Third, the design offers no way to force a
reset from software or a watchdog timer, a hung program
can only be recovered by unplugging the lamp.

\subsection{Component Tolerance}
\label{sec:tolerance}
Using a 5\% carbon-film resistor and a 20\% Y5V/Z5U ceramic
capacitor, both common, low-cost part choices
\[
\frac{\Delta t}{t} \approx \frac{\Delta R}{R} + \frac{\Delta C}{C}
= 0.05 + 0.20 = 0.25.
\]
Applied to the nominal $t_\text{release} = 18.1\,\text{ms}$,
the true release time falls in the range
\[
t_\text{release} \approx 13.6\,\text{ms to } 22.6\,\text{ms}.
\]

\begin{figure}[H]
\centering
\includegraphics[width=\columnwidth]{rc_tolerance.pdf}
\caption{Charging curves for the nominal, fastest-case, and
slowest-case component values.}
\label{fig:tolerance}
\end{figure}
\end{document}